% Minimal retained preamble for S3 SIGMA article cards taken from
% sources/sigma-2023-020-planar-orthogonal-polynomials-as-type-i-mult/pinned-original-source.tex
% (doi:10.3842/SIGMA.2023.020, SIGMA 19 (2023), CC BY 4.0).
%
% Omitted and why: the publisher \documentclass{sigma} block and the
% frontmatter commands (FirstPageHeading, ShortArticleName, ArticleName,
% Author, Address, Email, ArticleDates, Abstract, Keywords,
% Classification) - the sigma class is unavailable in the toolchain and
% the assembler supplies the article class; the dsfont package is
% dropped because no counted or context span of this card uses
% \mathds (unused-package disclosure per the three-pitfall broadcast,
% item ii); the TikZ settings block (cross style, tikzdeclarepattern)
% is dropped because no span of this card uses those styles.
% \samepage is a sigma-class spacing command the base LaTeX kernel
% does not define; author spans (source lines 123 and 831) contain it,
% so an empty stub is provided.
% The theorem declarations, equation numbering and paper macros below
% are copied verbatim from the pinned source. Note the counter name is
% Theorem (\newtheorem{Theorem}{Theorem}[section]); initial_counters
% keys use exactly these counter names.
\usepackage{amsmath,amssymb,amsthm}
\usepackage{mathtools}
\usepackage[mathscr]{euscript}
\providecommand{\qed}{\ensuremath{\square}}
\providecommand{\samepage}{}
\allowdisplaybreaks
\numberwithin{equation}{section}
\newtheorem{Theorem}{Theorem}[section]
\newtheorem*{Theorem*}{Theorem}
\newtheorem{Corollary}[Theorem]{Corollary}
\newtheorem{Lemma}[Theorem]{Lemma}
\newtheorem{Proposition}[Theorem]{Proposition}
\newtheorem*{rhproblem}{Riemann--Hilbert problem}
{ \theoremstyle{definition}
 \newtheorem{Definition}[Theorem]{Definition}
 \newtheorem{Note}[Theorem]{Note}
 \newtheorem{Example}[Theorem]{Example}
 \newtheorem{Remark}[Theorem]{Remark}
}
 \DeclareMathOperator{\Leb}{Leb}
 \DeclareMathOperator{\Span}{Span}
 \newcommand*{\Cset}{\mathbb{C}} % complex numbers
 \newcommand*{\Rset}{\mathbb{R}} % real numbers
 \newcommand*{\Nset}{\mathbb{N}} % natural numbers
